% Authoritative RouterBSI_Draft_Appendix.docx; zero-based direct body paragraph IDs.

% DOCX-P022-BEGIN
\section{Method}
% DOCX-P022-END

% DOCX-P023-BEGIN
\subsection{Predicate-Based Skill Routing}
% DOCX-P023-END

% DOCX-P025-BEGIN
Skills are selected according to whether the task aligns with the type of task each skill is intended to handle. This is determined by evaluating the skill's predicate statement for the given task. The predicate evaluates to true when the task aligns with the intended task type and to false otherwise.
% DOCX-P025-END

% DOCX-P026-BEGIN
The router first retrieves the top-k skills based on the retrieval score \(r(q,K_{s})\) between task \(q\) and the skill key \(K_{s}\). Given task \(q\), it then evaluates the predicate \(p_{sk}\) of each retrieved skill \(s_{k}\), where \(s_{k}\) ∈ \(S@k\), to determine whether it holds for \(q\): \(p_{sk}(q) = \{ True\ or\ False\}\). The router selects skills whose predicates evaluate to true: \(p_{sk}(q) = True\).
% DOCX-P026-END

% DOCX-P027-BEGIN
\subsection{MBR Update Selection with Bounded Replay}
% DOCX-P027-END

% DOCX-P028-BEGIN
Compare difference between \(P_{q}\)with \(G_{q}\) and identify the any of the following failures: required skill not selected, selected irrelevant skill. Skill \(s_{k}\) is selected only when predicate evaluates to true \(p_{s}\) given the task \(q\): \(p_{sk}(q) = True\). For each skill, the replay set contains previously seen relevant and irrelevant tasks:
% DOCX-P028-END

% DOCX-P029-BEGIN
\[\begin{gathered}\mathcal{R}_s = \{q_{rs}\}\cup\{q_{is}\},\\ |\{q_{rs}\}|\leq N,\qquad |\{q_{is}\}|\leq N\end{gathered}\]
% DOCX-P029-END

% DOCX-P030-BEGIN
For each failed skill \(s_{fk}\) in top \(k\), propose \(M_{p}\) candidate updates to each failed predicate \(p_{f}\) given the task \(q\), skill \(s_{sk}\), and \(q_{rs}\) relevant and \(q_{is}\) irrelevant tasks from previously seen tasks of the failed skill.
% DOCX-P030-END

% DOCX-P031-BEGIN
To reduce overly long and detailed predicates that may not generalize, we limit each proposed candidate \(m_{p} \in M_{p}\) must not exceeds a max length \(l\).
% DOCX-P031-END

% DOCX-P033-BEGIN
Replay each candidate \(m_{p} \in M_{p}\) on previously seen tasks where at most \(N\) are relevant and \(N\) are irrelevant for the skill \(s\).
% DOCX-P033-END

% DOCX-P034-BEGIN
After replay, use MBR to select an update \(m_{p}\) ∈ \(M_{p}\) or the original predicate \(p_{f}\) with the least risk score from the risk function. Our risk function is the average skill selection error on replayed tasks:
% DOCX-P034-END

% DOCX-P035-BEGIN
\[\ell(m_p,q,s)=|m_p(q)-\mathds{1}[s\in G_q]|\]
% DOCX-P035-END

% DOCX-P036-BEGIN
\[R_{\mathrm{replay}}(m_p\mid s)=\frac{1}{|\mathcal{R}_s|}\sum_{q\in\mathcal{R}_s}\ell(m_p,q,s)\]
% DOCX-P036-END

% DOCX-P037-BEGIN
\[m_p^*=\mathop{\arg\min}_{m_p\in M_p\cup\{p_f\}} R_{\mathrm{replay}}(m_p\mid s_{fk})\]
% DOCX-P037-END

% DOCX-P038-BEGIN
\[R_{\mathrm{replay}}(m_p^*\mid s_{fk})\leq R_{\mathrm{replay}}(p_f\mid s_{fk})\]
% DOCX-P038-END

% DOCX-P040-BEGIN
\subsection{Held-Out Validation and Recursive Improvement}
% DOCX-P040-END

% DOCX-P041-BEGIN
Evaluate the modified skill library \(S_{m}\) on a held-out validation set. Then, discard all selected updates \(M_{s}\) of modified skill \(s_{m} \in S_{m}\) that lead to a higher average selection error for the skill \(s_{m}\) on tasks where the skill \(s_{m}\) is included in top \(k\) for task \(t\). In other words, the average selection error is the difference between expected and predicated relevant/irrelevant label of the skill across tasks where the skill is in the top \(k\).
% DOCX-P041-END

% DOCX-P042-BEGIN
\[T_{v,s}=\{q\in T_v\mid s\in S@k\}\]
% DOCX-P042-END

% DOCX-P043-BEGIN
\[\begin{aligned}R_v(m_p\mid s)&=\frac{1}{|T_{v,s}|}\\&\quad\cdot\sum_{q\in T_{v,s}}|m_p(q)-\mathds{1}[s\in G_q]|\end{aligned}\]
% DOCX-P043-END

% DOCX-P044-BEGIN
\[\begin{aligned}&g(T_v,m_p^*,p_f)\\&\quad=\mathds{1}[R_v(m_p^*\mid s_{fk})\leq R_v(p_f\mid s_{fk})]\end{aligned}\]
% DOCX-P044-END

% DOCX-P045-BEGIN
\[p_f^{\mathrm{next}}=\begin{cases}m_p^*&\text{if }g(T_v,m_p^*,p_f)=1;\\p_f&\text{otherwise}\end{cases}\]
% DOCX-P045-END

% DOCX-P047-BEGIN
End the iteration by updating the skill library \emph{S} with the changes that passed the held-out validation gate. Repeat this self-improvement loop until i iterations: \(S^{(i+1)}=g(T_v,S_m)\)
% DOCX-P047-END
